\section{Discussion}
Idea:
\begin{itemize}
    \item Part 1:
    \begin{itemize}
        \item Hopper discovers forward locomotion BUT cannot reporoduce it in a stable way, therefore achieving at the very end of training very high rewards when it succeds in hopping and very low scores when it falls as it's not good ad hopping consistently.
        \item Discuss why spiky critic loss $\neq$ unstable training, training IS unstable but the spiky ciritic loss is not the proof
        \item Discuss insane actor loss stability
        \item STILL VERY AI GENERATED, I WAS VERY TIRED
        \item ADD IMPROVEMENTS
        \item Do I talk about the previous standing still and then collapsing hopper for AC? And that GAE + sigma floor fixed it
    \end{itemize}

    \item Part 2:
    \begin{itemize}
        \item Not sure, AI generated
    \end{itemize}
\end{itemize}

The results obtained in the first part demonstrate that the proposed Actor-Critic framework successfully learned a locomotion policy for the Hopper environment. Indeed, the Hopper development from balancing behavior to a stable forward locomotion, together with average rewards exceeding 2000 and occasional episodes above 3000, confirms that the agent discovered effective hopping. Nevertheless, training did not fully converge to a stable hopping policy: significant reward variability remained throughout the final stages of training, indicating that the learned policy was still too exploratory and stochastic, causing it to be unable to reproduce its best behaviours consistently.

This result represents a substantial improvement compared to earlier Actor-Critic implementations tested during development. Preliminary versions frequently converged to a degenerate local optimum in which the Hopper remained almost stationary for extended periods before eventually collapsing. Although such behaviour produced relatively stable trajectories, it prevented the discovery of effective locomotion and resulted in poor returns. The introduction of GAE, advantage normalization, a learned exploration variance with a sigma floor, and a triggered learning-rate decay mechanism significantly improved exploration and enabled the policy to escape this failure mode, eventually discovering sustained forward locomotion.

This instability seems to be confirmed by the critic-loss spikes, which, however should not necessarily be interpreted as training failure / instability. These spikes are not necessarily indicative of instability as reinforcement-learning targets continuously change because they depend on the current poliy, therefore when the actor discovers substantially better trajectories, previously learned value estimates become inaccurate, generating large temporal-difference errors and consequently large critic-loss spikes. This interpretation is supported by the fact that critic spikes often coincide with increases in reward, such as the discovery of transversal locomotion. In contrast, actor updates remained remarkably stable throughout training, as advantage normalization successfully bounded policy updates and prevented gradient explosions despite the highly non-stationary learning process.

Despite these improvements, the final policy still exhibits substantial variability and has not fully converged. Future work could investigate additional stabilization techniques like TRPO or PPO which exploits trust region and clipped policy updates \cite{schulman2017}.
Furthermore, more advanced actor-critic methods for continuous control, such as DDPG and TD3, could be explored. In particular, TD3 mitigates critic overestimation through double-critic architectures and delayed policy updates, potentially reducing value-estimation errors and improving training stability \cite{fujimoto2018}.


\textbf{(On this part I'm not sure but it's too long)}
The results of the second part show how SAC consistently outperformed PPO across all experiments. This behaviour is expected as SAC is off-policy and can repeatedly reuse past experience, substantially improving sample efficiency. Furthermore, the maximum-entropy objective encourages continued exploration and reduces the risk of premature convergence to suboptimal behaviours. PPO, on the other hand, is on-policy and discards collected trajectories after a limited number of updates. This lower sample efficiency appears to be a significant disadvantage and the experimental results confirm this expectation, with PPO failing to solve the task despite extensive hyperparameter tuning.

Domain Randomization proved effective in reducing the sim-to-real gap. By varying the cube mass during training, the policy could not memorize a specific strategy tailored to a single set of dynamics. Instead, it was forced to learn a more robust behaviour based on continuous observation and correction. Later on, the mass-sensitivity experiments confirmed the additional sturdiness introduced by domain randomization techniques as with the increment of the mass the vanilla SAC policy progressively deteriorated, while the UDR policy maintained a relatively stable performance across the entire tested range.

Contrary to the initial hypothesis, ADR did not outperform UDR. A plausible explanation could be that the manually selected UDR range [0.5,8.0] was a particularly nice pick as it already covered both source and target conditions effectively. In this setting, ADR had limited opportunity to discover a substantially better curriculum and may even have produced a less balanced sampling distribution during training. Therefore, the observed superiority of UDR should not be mistaked as generic but should be interpreted as task-specific.

Finally, the hyperparameter search suggests that SAC is relatively robust to parameter choices on this task. Neither the OFAT study nor the Bayesian optimization procedure produced a configuration that substantially outperformed the canonical SAC+UDR setup. This indicates that the default SAC configuration is already close to optimal for PandaPush and supports its use as the final deployed solution.
